% answers_short.tex —— 2006 年·普通高考·安徽卷理科 · 简短答案速查
% 由 tools/make_short_answers.py 自 answers.tex 的 \gksolhead/\gkanswer 提取，
% 并与 2006/2006_数学真题_安徽理_答案要点.md 逐题比对后生成，勿手改（改源 MD 后重跑）。

\begin{gknotice}
\textbf{使用说明}：本节只给最终答案便于快速核对，完整推导见解析版。答案为编写者独立推导，\textbf{非官方标准答案}。
\end{gknotice}

\begin{gknotice}
选择题答案中的 U 表示待核，未给出猜测答案，详情见解析版。
\end{gknotice}

\gkpart{选择题}
\begin{center}\begin{tabular}{ccccccccccccc}
\toprule
题号 & 1 & 2 & 3 & 4 & 5 & 6 & 7 & 8 & 9 & 10 & 11 & 12 \\
\midrule
答案 & A & B & U & B & C & C & A & B & A & B & D & C \\
\bottomrule
\end{tabular}\end{center}

\gkpart{填空题}
\begin{itemize}[leftmargin=2.6em]
  \item[\textbf{13.}] $1$。
  \item[\textbf{14.}] $\dfrac14(\vec b-\vec a)$。
  \item[\textbf{15.}] $-\dfrac15$。
  \item[\textbf{16.}] ①、③、④、⑤。
\end{itemize}

\gkpart{解答题}
\begin{itemize}[leftmargin=2.6em]
  \item[\textbf{17.}] (1) $-\dfrac13$；(2) $-\dfrac{5}{3\sqrt2}$
  \item[\textbf{18.}] (1) 当 $\xi=1,2,\ldots,9$ 时，其概率依次为 $\dfrac1{15},\dfrac1{15},\dfrac2{15},\dfrac2{15},\dfrac3{15},\dfrac2{15},\dfrac2{15},\dfrac1{15},\dfrac1{15}$；(2) $E\xi=5$
  \item[\textbf{19.}] (1) 见解析版；(2) 按包含 $A,D$ 的两个半平面，二面角为 $\arccos\left(-\dfrac3{\sqrt{105}}\right)$；若取两平面间的较小夹角，则为其补角
  \item[\textbf{20.}] (1) $f(0)=0$；(2) $f(x)=\begin{cases}kx,&x\geqslant0,\\hx,&x<0,\end{cases}$，其中 $k,h$ 为常数；(3) 在 $(0,1/k]$ 上递减，在 $[1/k,+\infty)$ 上递增；最小值为 $2$，在 $x=1/k$ 处取得，无最大值
  \item[\textbf{21.}] (1) $S_n=\dfrac{n^2}{n+1}$，递推关系为 $S_n=\dfrac{n^2}{n^2-1}S_{n-1}+\dfrac{n}{n+1}$（$n\geqslant2$）；(2) $T_n=\dfrac{p\left[1-(n+1)p^n+np^{n+1}\right]}{(1-p)^2}$（$p\ne1$）；$p=1$ 时 $T_n=\dfrac{n(n+1)}2$
  \item[\textbf{22.}] (1) $\lambda=e-\dfrac2e$，即 $e=\dfrac{\lambda+\sqrt{\lambda^2+8}}2$；由 $P$ 位于右支且在 $x$ 轴上方，还须 $e>\dfrac{1+\sqrt5}{2}$；(2) $x^2-\dfrac{y^2}{3}=1$
\end{itemize}
